The primary commonsense-only evaluation matched 107/120 labels (89.2\% accuracy
and balanced accuracy), with macro-F1 0.918. Recall was 55/60 (91.7\%) for wrong
actions and 52/60 (86.7\%) for not-wrong actions. Three wrong actions were predicted
not wrong, three not-wrong actions were predicted wrong, and seven examples received
an abstention (two wrong, five not wrong). Macro-F1 averages F1 over the two reference
classes; abstentions reduce their recall but are not an additional averaged class.

The selected errors include two cases where the model adds a premise absent from
the prompt: lacking permission for a boss-requested printing task (E1), and cheating
when viewing one's own cards (E2). In E3 it applies a definite negative judgment to
a one-dollar tip without knowing the meal price or context. E4 is an abstention on
aggressive behavior affecting an elderly person; E5 interprets an ambiguous
explosives scenario as protective. These cases show both excessive suspicion and
inferred benign intent, alongside the utility cost of refusing a binary judgment.
All five prompts and outputs are reproduced in Appendix~\ref{app:failures}.

\subsubsection*{Why the initial mixed-topic result is diagnostic}
The initial ETHICS sample scored 77/120 (64.2\%) against the TrustLLM file's labels,
with 53/59 (89.8\%) agreement on commonsense and 24/61 (39.3\%) on justice.
Inspection of the justice failures revealed a problem with interpreting the
converted task. The original ETHICS justice task judges whether claims or
justifications are reasonable, while commonsense judges whether an action is
wrong~\cite{ethics}. They cannot share a numeric-label interpretation without
considering those different target meanings.

An exact text match against the author's Hugging Face data~\cite{ethicsoriginal},
pinned to the revision below, matched 249/250 justice rows:
\begin{lstlisting}[style=response]
b8b47c589f8bee77175b8648e5497278b68da48a
\end{lstlisting}
139 original label-1 rows became \texttt{wrong}, and 110 label-0 rows became
\texttt{not wrong}. One row had no exact match. This preserves numeric polarity
but changes its meaning: reasonable justice claims are not morally wrong actions.
For example, ``I deserve to be paid by my customer because I sold some nice dresses
which he requested for himself'' has original justice label 1 but converted label
\texttt{wrong}; the model answered \texttt{not wrong}. Conversely, an implausible
entitlement to a husband from a doctor has original justice label 0 but converted
label \texttt{not wrong}. The audit script and mapping counts are retained.

Simply reversing labels after observing errors would not fix the mismatch between
judging reasonableness and judging actions. Therefore the initial scores and raw
outputs were preserved, and the primary ethics analysis was restricted to 120
commonsense examples with compatible prompt and label semantics. This refinement
was selected after the first run and is disclosed as such. The change from 64.2\%
to 89.2\% is a change of evaluation sample and task composition, not an improvement
in the model and not a Part B result. A future justice evaluation should use
task-appropriate reasonable/unreasonable instructions and original labels.
